\section{A simultaneous history selector with a direct iPRA transfer}
\label{sec:history-selector-proof}

Let $M$ be the pinned primary PA194 table in native numbering, with start
74 and terminal label 254.  The new machine $N=\mathrm{PA194SEL}$ keeps
the same 194 working labels and the same binary blank tape convention.
Its modification uses the bit written by the preceding transition to
remember which of two old controls has been entered.  The resulting
proof compares literal Turing-machine runs directly; it does not pass
through the numerical calculus or a proof-search scheduler.

\subsection{The complete finite certificate}

The following two pairs are used simultaneously:
\[
\begin{array}{c|c|c|c|c|c|c}
\text{pair}&l&r&\delta(l,0)&\delta(r,0)&u&v\\\hline
A&12&38&0R19&0R28&98&19\\
B&127&124&0R14&0R62&8&59
\end{array}
\]
For pair A, the two old read-one rows are both $1R20$; for pair B they
are both $1R138$.  The helper rows satisfy
\[
 \delta(98,0)=0R19,\quad \delta(19,0)=0R28,\quad
 \delta(8,0)=0R14,\quad \delta(59,0)=0R62.
\]
Every transition entering $l$ writes zero and moves right, and every
transition entering $r$ writes one and moves right.  These statements
are checked over the entire table.  Their complete predecessor lists are
\[
\begin{array}{c|l}
12&(2002,0),(2010,0),(2025,0),(2031,0)\\
38&(312,1)\\
127&(110,0)\\
124&(104,1).
\end{array}
\]
Thus, whenever an old run started at 74 is in $l$, the cell immediately
left of the head contains zero; when it is in $r$, that cell contains one.
The assertion holds for every starting tape.  It does not assert that an
arbitrary configuration initially placed in one of these four controls
already satisfies the required history condition.

Define $f(38)=12$, $f(124)=127$, and let $f$ fix every other label,
including 74 and 254.  Redirect every ordinary target $q$ to $f(q)$.
For each pair replace the merged control's read-zero row and the old
right control's two rows by
\[
 \delta_N(l,0)=0Lr,\qquad
 \delta_N(r,0)=0Ru,\qquad
 \delta_N(r,1)=1Rv.
\]
The read-one row of $l$ retains its old action with its target mapped by
$f$.  The former control $r$ is now a selector.  The old transitions
entering $r$ instead enter $l$, so they do not accidentally execute the
selector on an ordinary checkpoint entry.

Exactly eight literal rows change.  The helpers
$98,19,8,59$ are disjoint from the four modified control labels
$12,38,127,124$, and their relevant read-zero rows are unchanged by either
rewrite.  These facts are sufficient to check the two rewrites together;
there is no assumption that independent edits automatically compose.

\subsection{Whole-tape local simulation}

At a checkpoint replace only the old state by $f(q)$, leaving the head
and entire tape unchanged.  Each old step has the following exact new
simulation.
\begin{itemize}
\item Outside the four paired controls, one new transition reproduces
      the old write and move and reaches the mapped target.
\item At either control of a pair on read one, the common row gives the
      same one-step correspondence.
\item On read zero, $l$ moves left into selector $r$, without changing
      the zero just read.  The selector reads the history bit.  It writes
      that same bit back and moves right into $u$ if the bit is zero or
      $v$ if it is one.  The head is again over the original zero.
      The helper's read-zero row then performs the appropriate old
      transition.
\end{itemize}
The last case takes exactly three steps.  It accesses only the original
cell and its immediate left neighbor; neither bit changes.  The final
head is one cell to the right of the old head, and the final state is
the mapped old destination.  Every cell outside that two-cell window is
unaccessed.  The strict intermediate controls are a selector and a
helper, all working controls, so no premature halt occurs.

This gives eight nontrivial local cases, four for each pair.  Every case
quantifies the whole uninspected exterior.  The finite checker derives
the cases from the literal old and new tables and checks every remaining
row against target redirection.

\subsection{A quantifier-free arithmetical proof}

For explicit arithmetization, represent a finite-support tape relative
to its head by two natural numbers.  Let $L$ encode the cells strictly
to the left, nearest first in binary, and let $R$ encode the current
cell followed by the cells to the right, again nearest first.  A
configuration is $(q,L,R)$; the current symbol is $R\bmod2$, and the
history bit is $L\bmod2$.  If $R=2r+b$ and $L=2\ell+a$, a transition
writing $w$ has tape effects
\[
 \begin{array}{ll}
 \text{move right}:&(L,R)\mapsto(2L+w,r),\\
 \text{move left}: &(L,R)\mapsto(\ell,4r+2w+a).
 \end{array}
\]
These are primitive-recursive operations with elementary quantifier-free
defining identities.  An absolute head coordinate can be retained as
an additional integer code; each simulated block makes the same net
move, so this does not change the argument.

Make terminal label 254 absorbing in both arithmetic step functions.
Write $c_n$ for the old configuration after $n$ steps, and let
$F(q,L,R)=(f(q),L,R)$.  The entry invariant is the quantifier-free formula
\[
 J(q,L,R):\quad
 (q=12\lor q=127\Rightarrow L\bmod2=0)\ \wedge\
 (q=38\lor q=124\Rightarrow L\bmod2=1).
\]
The initial state 74 satisfies $J$.  Every transition into a paired
control writes its specified bit and moves right, so the complete
predecessor check proves that one old step preserves $J$, for arbitrary
values of $L,R$.  Quantifier-free induction gives $J(c_n)$ for every $n$.

Define the primitive-recursive clock
\[
 \tau(0)=0,\qquad
 \tau(n+1)=\tau(n)+d(c_n),
\]
where $d(c)=3$ when $c$ is in one of the four paired controls and reads
zero, and $d(c)=1$ otherwise.  The absorbing halt case has duration one.
The local certificate and quantifier-free induction establish
\[
 \operatorname{run}_N(\tau(n))=F(c_n),\qquad
 n\le\tau(n)\le3n.                                      \tag{1}
\]
The induction formula is an equality between explicitly
primitive-recursive configuration codes, conjoined with $J(c_n)$ and the
two clock inequalities.  It has no unbounded quantifier over tapes,
stores, proofs, or future runs.

Let $h_M(t)$ and $h_N(t)$ mean that the respective absorbing runs are in
the halt state at time $t$.  From (1), if $h_M(t)$ then $N$ is halted at
$\tau(t)\le3t$ and remains halted thereafter.  Thus iPRA proves
\[
 \forall t\,[h_M(t)\to h_N(3t)].                          \tag{2}
\]

For the converse, suppose $h_N(T)$.  Since the halt state is absorbing
and $\tau(T)\ge T$, the new machine is also halted at $\tau(T)$.
Equation (1) then says that $F(c_T)$ is halted.  The map $f$ sends no
working state to 254, so $c_T$ itself is halted.  Hence
\[
 \forall T\,[h_N(T)\to h_M(T)].                          \tag{3}
\]
Halting persistence is itself a quantifier-free induction on the number
of additional absorbing steps.  No phase reconstruction or search for
the last checkpoint is needed.  Positivity of the block durations,
which gives $\tau(T)\ge T$, supplies the essential bound.

Equations (2)--(3) are explicit halting-witness translations.  Pure
intuitionistic reasoning now gives the direct relative theorem
\[
 \boxed{\mathsf{iPRA}\vdash H_M\leftrightarrow H_N,\qquad
        \mathsf{iPRA}\vdash\neg H_M\leftrightarrow\neg H_N.}
\]
This is a mathematical proof of the relative theorem; the accompanying
Python checker checks its complete finite table premises and is not a
new Lean kernel compilation.  No soundness principle, consistency
assumption, reflection principle, or completed PA-calibration theorem is
needed for the relative result.

Consequently an established primary-machine theorem
$\mathsf{iPRA}\vdash\operatorname{Con}(\mathsf{PA})
\leftrightarrow\neg H_M$ transfers to $N$ by intuitionistic
composition.  The current repository's missing absolute calibration
does not disappear: it is inherited unchanged, rather than replaced by
a new command-calculus correctness obligation.

\subsection{Blank initialization as an independent concrete check}

The old blank startup reaches its source frame after 58,154 transitions.
During that prefix the old run takes a read-zero transition in pair A
three times and in pair B twenty-four times.  Each such visit gains two
steps and all others gain none, so the new startup length is exactly
\[
 58154+2(3+24)=58208.
\]
Independent literal simulation confirms that both endpoints have state
84, head coordinate $-2$, and the identical complete finite-support tape.
This concrete check binds the reported startup to the actual files.
The all-time equivalence follows from the local proof and induction
above, not from extrapolating the startup computation.
